\chapter{Method}
\label{chap:method}
\section{Problem Formulation}

I frame the task as imitation learning from recorded gameplay demonstrations
\cite{Bain1995AFF}. At each timestep $t$, the agent observes an RGB frame
$s_t$ of resolution $360 \times 640 \times 3$, paired with a natural-language
task prompt $p$ describing the goal (e.g.\ \enquote{collect dirt}). The agent
must predict an action $a_t \in \{0,1\}^{21} \times \mathbb{R}^2$, comprising
21 binary key presses and 2 continuous camera magnitudes $(c_x, c_y)$. The
training signal is a dataset of demonstrations
\begin{equation}
    \mathcal{D} = \{(s_t,\, p,\, a_t)\}_{t=1}^{N},
\end{equation}
where each frame is paired with the action the demonstrator executed at that
timestep. No reward signal is available; the learner can only infer
$\hat{\pi}$ from the state-action pairs in $\mathcal{D}$. I therefore adopt
Behavioural Cloning (BC) as my learning algorithm.

The policy does not consume visual history: only the current frame $s_t$ is
encoded by the backbone. To partially compensate, the head is conditioned on
a short window of past actions and predicts a short chunk of future actions,
as described in Sections~\ref{past_action} and~\ref{chunking}.

Finally, I note that the absence of expert querying during training can lead
to \textit{compounding errors}: once the agent deviates from states seen in
$\mathcal{D}$, it finds itself in an out-of-distribution region from which it
cannot recover \cite{ross2011}. This offline-to-online gap is a fundamental
property of plain BC, where errors scale as $\mathcal{O}(\varepsilon T^2)$
with horizon $T$, and motivates both the action-chunking design choice in
Section~\ref{chunking} and the in-environment rollout evaluation in
Section~\ref{sec:rollout}.

\section{Architecture} \label{sec:architecture}

\subsection{Frozen Vision-Language Backbone} \label{sec:backbone}
Freezing strong pretrained representations and training only a lightweight
head has been shown to match or exceed full fine-tuning accuracy while
remaining computationally
tractable~\cite{kumar2022finetuningdistortpretrainedfeatures}, a paradigm
established in vision-language modelling by
BLIP-2~\cite{li2023blip2bootstrappinglanguageimagepretraining}. I adopt
LLaVA-1.5-7B~\cite{liu2024improvedbaselinesvisualinstruction}, which pairs a
CLIP-ViT-L-336px vision encoder with a Vicuna-7B language decoder, as the
primary frozen backbone. Each state $(s_t, p)$ is passed jointly through
\texttt{LlavaProcessor} and the autoregressive decoder, where the single
\texttt{<image>} placeholder is expanded to 576 visual tokens that are
processed alongside the prompt tokens through all decoder layers. I read the
post-normalisation hidden states of the final layer via a forward hook on the
decoder's terminal RMSNorm.

Because the LLaMA decoder exposes no CLS token, a fixed-size summary must be
pooled from the token sequence. A naive mean over the full joint sequence is
dominated by the roughly 576 visual tokens and washes out the handful of
prompt tokens. I therefore pool the image-role and text-role tokens
\emph{separately} using their attention masks and concatenate the two
4096-dimensional means:
\begin{equation}
    \phi_{\mathrm{LLaVA}}(s_t, p) = \bigl[\, \overline{h}^{(L)}_{\mathcal{I}}\,;\;
    \overline{h}^{(L)}_{\mathcal{T}} \,\bigr] \in \mathbb{R}^{8192},
\end{equation}
where $\overline{h}^{(L)}_{\mathcal{I}}$ and $\overline{h}^{(L)}_{\mathcal{T}}$
are the masked means of the final-layer hidden states over the image-role
token set $\mathcal{I}$ and the text-role token set $\mathcal{T}$, each of
dimension 4096. Split-pooling keeps the language signal as an explicit,
undiluted component of the representation, mirroring the concatenated layout
of the contrastive baseline below. All 7B backbone parameters are frozen
throughout training, so only the action head defined in
Section~\ref{action_head} receives gradient updates.

\subsection{Contrastive Vision-Language Baseline} \label{sec:clip_baseline}
To isolate the contribution of LLaVA's joint vision-language reasoning from
its raw visual features, I include a frozen CLIP-ViT-L/14
baseline~\cite{radford2021clip} as a comparison backbone. Whereas LLaVA fuses
image and text inside a single autoregressive decoder, CLIP encodes the two
modalities with separate encoders that have only been aligned by contrastive
pretraining. The image encoder produces a 768-dimensional pooled image
embedding $\phi_v(s_t)$ and the text encoder produces a 768-dimensional
pooled text embedding $\phi_t(p)$. I concatenate the two to form the joint
representation
\begin{equation}
    \phi_{\mathrm{CLIP}}(s_t, p) = \bigl[\,\phi_v(s_t);\; \phi_t(p)\,\bigr] \in \mathbb{R}^{1536}.
\end{equation}
When the prompt is ablated (Section~\ref{ablation}), the text branch is
replaced with a zero vector of the same dimensionality rather than removed,
so the head input shape is held constant across language conditions and only
the backbone choice is varied. Because the CLIP image encoder is separate
from the text encoder, its visual features never attend to the prompt; this
is the key respect in which the no-prompt condition differs between the two
backbones (Section~\ref{ablation}).

\subsection{Action Space} \label{action_space}
The MineRL environment exposes a hybrid action space inherited from
VPT~\cite{baker2022vpt}: 21 binary keys (movement, attack, use, hotbar
selection, inventory toggle, ESC) and 2 continuous camera magnitudes
$(c_x, c_y) \in \mathbb{R}^2$ controlling view pitch and yaw. Predicting the
21 binary keys jointly as a single categorical distribution over all $2^{21}$
combinations is intractable; following the action-branching paradigm of
Tavakoli et al.~\cite{tavakoli2018actionbranching}, I instead predict each
key independently via its own sigmoid output, which scales linearly rather
than combinatorially in the number of keys.

The continuous camera axes are handled separately. Rather than regressing
$(c_x, c_y)$ directly, I discretise each axis into 11 mu-law-spaced bins
following the VPT convention~\cite{baker2022vpt}, allowing the policy to
represent the bimodal \enquote{stay still vs.\ turn} structure of human
demonstrations that a regression loss would collapse to the mean. Each axis
is therefore predicted as a categorical distribution over 11 bins, decoded
back to degrees at rollout time via the mu-law inverse. This yields the
per-step output dimensionality $D = 21 + 2 \times 11 = 43$.

\subsection{Past-Action Conditioning} \label{past_action}
A single frame is often insufficient to disambiguate the current action under
behavioural cloning. Tasks such as \texttt{collect\_dirt} and
\texttt{chop\_a\_tree} are dominated by long, repetitive action runs:
attack sequences of roughly 30 ticks while breaking a block, forward runs of
50 to 200 ticks while traversing the world, and small camera adjustments when
tracking a stationary target. Within such a run, consecutive frames are
visually near-identical, yet the correct action is overwhelmingly to repeat
the previous one. The conditional probability $\Pr(a_t = 1 \mid a_{t-1} = 1)$
is much higher than the marginal $\Pr(a_t = 1)$ for these sticky actions, but
single-frame BC only learns the marginal.

I therefore condition the action head on the last $K = 8$ executed actions.
Each past action is encoded as a $D = 43$ dimensional vector matching the
per-step output layout: the 21 binary keys as $0/1$ indicators (several of
which may be active at once) and each of the two camera axes as a one-hot
over its 11 mu-law bins. The resulting $K \times D$ matrix is flattened and
concatenated to the pooled backbone feature $\phi(s_t, p)$ before the head.
The window is zero-padded at trajectory start and right-aligned so the most
recent action occupies the last slot. This signal is intentionally injected
at the head input rather than appended to the text prompt. The head receives
the $K \cdot D = 344$ dimensional past-action vector as a dedicated
full-strength input to its trainable weights, whereas the same information
expressed as prompt text would have to be encoded and summarised by the
frozen backbone, competing with the task prompt inside the pooled text
representation and reaching the head only as a diluted, fixed-width fragment.

\subsection{Action Chunking} \label{chunking}
To mitigate compounding error and exploit short-horizon predictability, the
head predicts a chunk of $N = 8$ future actions per forward pass rather than
a single next action. Concretely, the head output is reshaped from a flat
vector of size $N \cdot D$ into an $(N, D)$ tensor, with each row interpreted
as the logits for one future timestep. During training, the loss averages
over the chunk axis so every predicted step contributes equally. At rollout
time the policy executes only the first step of each chunk before re-planning
on the next frame; the remaining steps serve only as auxiliary supervision
that encourages the head to internalise short-horizon dynamics.

\subsection{Trainable Action Head} \label{action_head}
The action head is a two-layer MLP with ReLU activation that maps the
concatenated input
\begin{equation}
    x_t = \bigl[\,\phi(s_t, p);\; \mathrm{vec}(a_{t-K:t-1})\,\bigr]
    \in \mathbb{R}^{F + K \cdot D},
\end{equation}
where $F$ is the backbone feature dimensionality ($F = 8192$ for LLaVA,
$F = 1536$ for CLIP), to a vector of chunked action logits
\begin{equation}
    \pi_\theta(s_t, p, a_{t-K:t-1}) = W_2 \, \mathrm{ReLU}(W_1 \, x_t + b_1) + b_2,
\end{equation}
which is reshaped to $\mathbb{R}^{N \times D}$. With $K = 8$ past actions,
$N = 8$ chunk steps, hidden width $H = 2048$, and per-step dimension
$D = 43$, the head has $W_1 \in \mathbb{R}^{(F + K D) \times H}$ and
$W_2 \in \mathbb{R}^{H \times (N D)}$, totalling approximately $18.2$M
parameters for the LLaVA variant and $4.6$M parameters for the CLIP variant,
both negligible against the $7$B and $0.43$B frozen backbones respectively.
% Confirmed against the four reported checkpoints (verification_paper.md,
% 2026-06-23, configs dumped from model_best.pt): hidden_dim=2048 for all four
% cells; head W1 shapes (2048, 8536) LLaVA / (2048, 1880) CLIP; head parameter
% counts exactly 18,188,632 (LLaVA) and 4,557,144 (CLIP). The trained head is the
% cached HeadOnlyAgent, hidden width held constant across both backbones, not the
% embed-dim-fixed (768) head in frozen_vision_baseline.py.
The action head is held in \texttt{fp32}, and pooled backbone
features are cast to \texttt{fp32} at the head input. During feature
extraction the frozen LLaVA decoder computes in \texttt{bf16}
(Section~\ref{caching}); the CLIP baseline runs in its native \texttt{fp32}.

\section{Training Objective}

\subsection{Loss Function} \label{sec:loss}
The Minecraft action space is fundamentally multi-label: a human player
routinely combines actions, for example holding \texttt{forward} while
pressing \texttt{attack} and adjusting the camera. A softmax cross-entropy
loss would force the policy to commit to a single most-likely action,
collapsing this combinatorial structure to a one-hot prediction. Following
the action-branching rationale of Tavakoli et
al.~\cite{tavakoli2018actionbranching}, I treat the 21 binary keys at each
chunk step as independent Bernoulli outputs trained with binary
cross-entropy. Let $z_{t,n} \in \mathbb{R}^D$ denote the predicted logits
for chunk step $n$ at training timestep $t$, with binary slice
$z^{\mathrm{bin}}_{t,n} \in \mathbb{R}^{21}$ and target
$a^{\mathrm{bin}}_{t,n} \in \{0,1\}^{21}$. The binary loss is
\begin{equation}
    \mathcal{L}_{\mathrm{bin}}(\theta) = -\frac{1}{N_{\mathrm{tot}} \cdot N \cdot 21}
    \sum_{t=1}^{N_{\mathrm{tot}}} \sum_{n=1}^{N} \sum_{i=1}^{21}
    \Bigl[ a^{\mathrm{bin}}_{t,n,i} \log \sigma(z^{\mathrm{bin}}_{t,n,i})
         + (1 - a^{\mathrm{bin}}_{t,n,i}) \log(1 - \sigma(z^{\mathrm{bin}}_{t,n,i})) \Bigr],
\end{equation}
where $N_{\mathrm{tot}}$ is the number of training samples, $N$ is the chunk
size, and $\sigma$ is the sigmoid function.

The two camera axes are categorical (Section~\ref{action_space}) and trained
with standard cross-entropy over their 11 mu-law bins. Let
$\tilde{a}^{(x)}_{t,n}, \tilde{a}^{(y)}_{t,n} \in \{1, \dots, 11\}$ denote
the target bin indices for camera_x and camera_y at chunk step $n$, and let
$z^{(x)}_{t,n}, z^{(y)}_{t,n} \in \mathbb{R}^{11}$ denote the corresponding
logit slices. The camera loss is the average over both axes and all chunk
steps:
\begin{equation}
    \mathcal{L}_{\mathrm{cam}}(\theta) = \frac{1}{2 N_{\mathrm{tot}} N}
    \sum_{t=1}^{N_{\mathrm{tot}}} \sum_{n=1}^{N}
    \Bigl[ \mathrm{CE}\bigl(z^{(x)}_{t,n}, \tilde{a}^{(x)}_{t,n}\bigr)
         + \mathrm{CE}\bigl(z^{(y)}_{t,n}, \tilde{a}^{(y)}_{t,n}\bigr) \Bigr].
\end{equation}

The full training objective is the unweighted sum:
\begin{equation}
    \mathcal{L}(\theta) = \mathcal{L}_{\mathrm{bin}}(\theta) + \mathcal{L}_{\mathrm{cam}}(\theta).
\end{equation}
Averaging the two camera cross-entropies internally keeps
$\mathcal{L}_{\mathrm{cam}}$ on roughly the same scale as
$\mathcal{L}_{\mathrm{bin}}$ from the start of training, removing the need
for a hand-tuned weighting coefficient.\footnote{The implementation writes
this as $\mathcal{L}_{\mathrm{bin}} + 0.5\,(\mathrm{CE}_x + \mathrm{CE}_y)$,
which is algebraically identical to summing $\mathcal{L}_{\mathrm{bin}}$
with the mean of the two camera cross-entropies as expressed here.}

\subsection{Optimisation}
Only the action head parameters $\theta$ receive gradient updates; the
frozen backbone parameters remain unchanged throughout. I optimise with Adam
(learning rate $\eta = 10^{-3}$, default $\beta = (0.9, 0.999)$, no weight
decay) for 10 epochs at a batch size of 256 under a constant learning-rate
schedule, with no gradient clipping or warm-up. Because the head is trained
against pre-extracted features (Section~\ref{caching}), no mixed-precision
autocast is applied at this stage: the cached \texttt{fp16} features are cast
to \texttt{fp32} and the head trains entirely in \texttt{fp32}. I select the
final checkpoint by the epoch with the highest validation movement-$F_1$, the
mean $F_1$ over the locomotion keys (back, forward, jump, left, right,
sprint), since the attack key saturates early and is not discriminative for
model selection.

To separate the effect of the ablated factors from training noise, every cell is
trained at \textbf{three independent seeds}. The seed varies the head's random
initialisation and the data-loader shuffle order; the backbone, the
train/validation/test partition (Section~\ref{sec:split}), and the optimiser
settings above are held fixed across seeds. Because the two-stage design makes
head training cheap (Section~\ref{caching}), this triples only the inexpensive
stage and leaves the one-time feature cache untouched. Throughout
Chapter~\ref{chap:experiments} every reported quantity is the mean over the three
seeds, with the across-seed standard deviation as its uncertainty, so an effect
counts only if it survives retraining rather than holding for one lucky seed.

\subsection{Two-Stage Training with Cached Features} \label{caching}
Running the backbone end-to-end for every training step is prohibitively
expensive: a LLaVA forward pass costs on the order of $15$~ms per sample,
and running the full ablation end-to-end would cost roughly $53$ times the
cached pipeline. Because the backbone is frozen, however, its output
$\phi(s_t, p)$ depends only on $(s_t, p)$, both of which are fixed across all
head-training runs for a given dataset and language condition. I therefore
exploit this invariance with a two-stage training pipeline.

\paragraph{Stage 1: feature caching.} For each backbone-and-language cell I
precompute $\phi(s_t, p)$ for every frame in $\mathcal{D}$ (at the production
frame-stride of 4) and store the resulting embeddings as an \texttt{fp16}
memory-mapped array, accompanied by a JSON sidecar that pins the sample
ordering and the cache configuration. This yields four caches in total, and
the backbone runs exactly once per cache. The no-prompt condition is not
produced by feeding an empty prompt: for LLaVA the full prompt is always
passed through the decoder and only the pooled text component is zeroed after
pooling, while for CLIP the text branch is replaced by a zero vector
(Section~\ref{ablation}). Building a cache costs roughly $6$~hours for LLaVA
and roughly $1.5$~hours for the CLIP baseline.

\paragraph{Stage 2: head training.} The action head is then trained against
the cached features at MLP speed, orders of magnitude faster than an
end-to-end forward pass. Crucially, the cache is invariant to the past-action
window size $K$ and the chunk size $N$, since both affect only the head, not
the backbone. The same cached features therefore serve every head-architecture
variant explored in the ablation. With the cache in place, head training
itself costs on the order of $0.5$ to $2$~hours per cell, leaving the
one-time cache build as the dominant cost of the whole ablation.

\section{Data Pipeline}

\subsection{Dataset}
Demonstrations come from recorded gameplay on two tasks,
\texttt{chop\_a\_tree} and \texttt{collect\_dirt}. The trajectories are not
human recordings: they are real MineRL point-of-view renders produced by a
STEVE-1 / VPT agent~\cite{lifshitz2023steve1, baker2022vpt} driving the live
game engine through the MineDreamer \texttt{play} pipeline, so the frames are
genuine engine output rather than generative-model imagery.
% Provenance confirmed (verification_paper.md): real MineRL engine renders from a
% MineDreamer play run driven by STEVE-1/VPT (in_model vpt/2x.model, in_weights
% steve1.weights), not diffusion. NOTE: dbstmpl.bib is still the stock LMU
% template, so no MineDreamer key exists and ALL cite keys here are currently
% undefined; the bibliography must be populated before the chapter compiles
% (see open-items block).
Each trajectory consists of a $20$~FPS, $360 \times 640$ MP4 video paired
with a per-frame action log in JSON format and a single task-prompt string in
the trajectory metadata; frames and actions are aligned one-to-one. Because
the generating policy is VPT-based, the recorded camera magnitudes already
coincide with the mu-law bin centres of Section~\ref{action_space}, so
discretisation to the 11-bin representation is lossless.

In total the combined dataset comprises $2{,}081$ trajectories --- $1{,}065$ for
\texttt{chop\_a\_tree} and $1{,}016$ for \texttt{collect\_dirt} --- of $3{,}000$
frames each, for $6{,}243{,}000$ frames or roughly $87$~hours of $20$~FPS play,
split almost evenly between the two tasks ($51.2\%$ chop, $48.8\%$ dirt). The
near-even task balance matters for the language ablation of
Section~\ref{ablation}: because the prompt is tied to whichever task a trajectory
belongs to, a balanced corpus ensures that a model which conditions on the prompt
cannot simply be exploiting a frequency prior over the two prompt strings. At the
production frame-stride of $4$ this yields the $1{,}560{,}750$ cached training
samples per cell used in Section~\ref{caching}.

\paragraph{Channel ordering.} The training videos were written with the red
and blue channels swapped, an artefact of saving frames through OpenCV, which
expects BGR rather than RGB byte order, so the policy is trained on
R/B-swapped frames. The same swap is therefore applied to the live
observation at rollout time, making it the only systematic difference between
the training and evaluation frame distributions.

\subsection{Frame Loading}
Whenever the backbone consumes raw frames, that is during the one-time cache
build and the end-to-end sanity path, frames are decoded on demand from the
source MP4 using \texttt{decord.VideoReader} inside the PyTorch
\texttt{Dataset} rather than pre-extracted to disk as PNG or HDF5. Each
DataLoader worker maintains its own LRU cache of recently opened video
readers (default size $64$), so the amortised random-access latency per frame
is in the $3$ to $5$~ms range, well below the backbone forward pass and thus
invisible in end-to-end timing. This avoids both the disk-space blow-up and
the inode-explosion problems of materialised frame extraction. Head training
reads pre-pooled features from the cache (Section~\ref{caching}) and does not
decode frames.

\subsection{Data Split} \label{sec:split}
The dataset is split into training, validation, and test sets with approximate
proportions $80\% / 10\% / 10\%$, using a fixed random seed (42) for
reproducibility. The
split is performed at the \emph{trajectory} level rather than the frame level:
every frame of a given trajectory is assigned to exactly one of the three
sets. This is deliberate. Adjacent frames in Minecraft gameplay are visually
near-identical, so a frame-level split would scatter near-duplicates of test
frames into the training set and inflate held-out metrics through leakage.
Trajectory-level splitting holds the training and test trajectories strictly
disjoint, so the offline metrics of Section~\ref{sec:rollout} measure
generalisation to unseen trajectories rather than memorisation of neighbouring
frames. I additionally rely on the in-environment rollout evaluation of
Section~\ref{sec:rollout} as a still stronger test: rollout episodes execute
the policy in freshly generated worlds that share no frames with any
trajectory in $\mathcal{D}$, probing out-of-distribution behaviour beyond what
any held-out split of the recorded data can capture.

\section{Rollout Evaluation} \label{sec:rollout}

\subsection{Inference Loop}
At evaluation time the trained head is paired with its frozen backbone and
deployed in the MineRL environment $\mathcal{E}$. Each episode maintains a
fixed-length deque of the last $K = 8$ executed actions, zero-initialised at
$t = 0$. At each timestep the agent encodes the current frame and prompt to
obtain $\phi(s_t, p)$, concatenates the past-action buffer, and forwards
through the head to produce a chunked logit tensor of shape $(1, N, D)$.
Only the first chunk step is decoded into an environment action. I decode
\emph{stochastically}: each binary key is sampled from a Bernoulli with
probability $\sigma(z / \tau)$, and each camera axis is sampled from a
categorical over its 11 bins via $\mathrm{softmax}(z / \tau_{\mathrm{cam}})$
before the mu-law inverse maps the chosen bin back to degrees. I use
$\tau = 1.0$ for the binary keys and a higher camera temperature
$\tau_{\mathrm{cam}} = 2.0$, which flattens the dominant zero-magnitude bin so
the agent looks around rather than fixating on the most frequent
\enquote{no turn} action; a greedy variant (binary threshold at $0.5$, camera
\texttt{argmax}) is available but tends to collapse the camera onto the null
bin. The resulting action is merged into the environment's no-op template
(which fills in components such as \texttt{pickItem} and \texttt{swapHands}
that the policy does not predict), executed in $\mathcal{E}$, appended to the
past-action buffer as a one-hot vector, and the cycle repeats. The policy
therefore re-plans on every tick; the remaining $N - 1$ chunk steps act only
as auxiliary supervision during training (Section~\ref{chunking}) and are
discarded at rollout.

\subsection{Evaluation Metrics}
I report two complementary classes of metric, and treat the online rollout
statistics as primary.

Offline frame-level metrics are computed against the held-out test split of
Section~\ref{sec:split}: per binary key I report precision, recall, and
$F_1$, and for each camera axis I report top-1 bin accuracy and mean absolute
error in degrees after decoding. These provide a cheap signal during
development but are known to correlate weakly with downstream task success
\cite{kanervisto2020}.

Online rollout metrics are evaluated across a fixed number of episodes per
cell. The primary instrument is the \emph{per-action firing rate}, the
fraction of steps at which each binary key fires, measured per condition; I
compare firing rates across conditions using Wilson confidence intervals, a
two-proportion $z$-test with multiple-comparison correction, and an
effect-size threshold, so a difference must be both statistically and
practically meaningful to count. I favour firing-rate statistics over episode
reward because the in-environment reward signal is both systematically broken on
the chop task and a small sample at the episode counts that are computationally
affordable: reward must be recomputed from inventory deltas because the engine's
item-collection reward handler watches the generic item name \texttt{log} while
chopping yields per-wood names such as \texttt{oak\_log} and \texttt{birch\_log},
so it never increments for the chop task in this stack, and even on the working
dirt task ten episodes is a small sample over ten distinct world seeds. The
environment is itself seed-deterministic --- an identical seed reproduces a
rollout exactly --- so the per-episode spread is sampling variance over worlds,
not irreproducibility. As a reward-independent measure of task progress I track the per-item
\emph{peak inventory}, the running maximum count of the goal item (a log or a
dirt block) over an episode, alongside the recomputed episode reward and
success rate reported for completeness. As a sanity floor I record the
action-frequency distribution and compare it against a random-policy baseline
run in the same environments, confirming that the trained agents diverge from
chance (Chapter~\ref{chap:experiments}).
% [RESOLVED 2026-06-24] random-policy baseline now logged at
% evaluations/paper/random_baseline/ (chop+dirt, 10 episodes each, seed 0;
% binary key flip p=0.1, camera ~U(-5,5)). Floor ~10%/key, 0 logs / 0 dirt
% collected; reported in Chapter~\ref{chap:experiments} (Sec. Experimental Setup,
% Fig. action profile). The stronger "I ... confirm ... diverge from chance"
% wording is therefore restored above.

\section{Ablation Design} \label{ablation}

To disentangle the contribution of the joint vision-language backbone from
the contribution of the language prompt itself, I run a $2 \times 2$
factorial ablation over backbone $\in \{\text{LLaVA}, \text{CLIP}\}$ and
language condition $\in \{\text{prompt}, \text{no prompt}\}$, giving four
head-training cells. All four cells are trained on the \emph{combined}
two-task dataset (\texttt{chop\_a\_tree} and \texttt{collect\_dirt} pooled
together), not on either task in isolation. Training on the combined set is
what makes language conditioning a non-trivial question at all: on a single
task the prompt is constant, so for the contrastive baseline it would fold
into a fixed head bias and the prompt-versus-no-prompt comparison would be
vacuous.

Across all four cells the temporal recipe (past-action conditioning with
$K = 8$ and action chunking with $N = 8$), the head hidden width, the
optimiser, the frame stride, the trajectory-level split, and the
model-selection criterion are held identical; only the backbone identity and
the language condition vary, so that neither backbone is handicapped by the
comparison.

The four cells are:
\begin{itemize}
    \item \textbf{Exp 1}, LLaVA with prompt: the full model
          (Section~\ref{sec:backbone}).
    \item \textbf{Exp 2}, CLIP with prompt: the contrastive baseline with
          access to the same language signal.
    \item \textbf{Exp 3}, LLaVA without prompt: the prompt is still passed
          through the decoder, but the pooled text component is zeroed after
          pooling (see below).
    \item \textbf{Exp 4}, CLIP without prompt: the text branch is replaced by
          a zero embedding while the head input shape is held unchanged
          (Section~\ref{sec:clip_baseline}).
\end{itemize}

\paragraph{What the language ablation removes, per backbone.} The no-prompt
condition is not implemented identically across the two backbones, and the
asymmetry matters for interpretation. For CLIP, the image and text encoders
are separate, so zeroing the text branch removes language from the
representation entirely: the image features never attended to the prompt. For
LLaVA, the prompt is always passed through the decoder and only the pooled
text component is zeroed afterwards (Section~\ref{sec:backbone}). Because
LLaVA fuses the modalities inside the decoder, the pooled \emph{image}
component has already attended to the prompt and so remains prompt-conditioned
even in the no-prompt cell. The LLaVA language ablation therefore removes the
explicit text vector but not the prompt's influence on the visual
representation, whereas the CLIP ablation removes language completely.

\paragraph{What each comparison isolates.} Exp 1 versus Exp 2 holds the
language signal fixed and varies only the backbone, isolating the value of
joint vision-language fusion against a contrastive encoder; this is the
cleanest contrast the design supports. Exp 1 versus Exp 3, and Exp 2 versus
Exp 4, each measure the marginal value of the prompt within a single backbone,
and Exp 3 versus Exp 4 compares the two backbones as pure visual feature
extractors. One caveat must be stated plainly: because all cells are trained
on the combined dataset, the prompt also carries the identity of the task.
The within-backbone prompt comparisons therefore measure the total value of
the prompt, including its role as a task-disambiguation signal, and cannot on
their own attribute any gain specifically to language reasoning. This is a
deliberate consequence of keeping the design a tidy $2 \times 2$: a fifth cell
supplying task identity through a non-linguistic one-hot signal would have
separated the two effects, but I do not include it. The overarching question
the design poses is whether a frozen joint vision-language model earns its
keep over a frozen contrastive one when both are used only as feature
extractors for behavioural cloning, and whether either backbone makes
meaningful use of the language pathway under this training regime.

\section{Implementation Details}
The system is implemented in PyTorch on Python 3.12. Feature caching and head
training run on a rented NVIDIA RTX 5090, using \texttt{bf16} compute with
% Hardware/stack confirmed (verification_paper.md): RTX 5090 (Blackwell, sm_120),
% torch 2.7 from the cu128 index (CUDA 12.8), bf16 + SDPA, no FlashAttention-2.
% The A5000/Abaki path is deprecated (QoS revoked 2026-06-06); the four reported
% cells are stride-4 LLaVA/CLIP runs consistent with the 5090 build.
scaled-dot-product attention for the LLaVA backbone (the Blackwell
architecture has no FlashAttention-2 kernel). Rollout evaluation runs in a
Dockerised MineRL environment. The full codebase, including the cache
pipeline, training loop, and rollout harness, is available at
\url{https://github.com/DACN-1/MC-AI}.

% =====================================================================
% OPEN ITEMS (Method) -- status after verification_paper.md (2026-06-23),
% which dumped the four reported checkpoint configs directly from model_best.pt.
% ---------------------------------------------------------------------
% RESOLVED -- confirmed against source AND the four reported checkpoints
% {llava,clip}_combined_{lang,nolang}_stride4_tsplit[_base]:
%  * hidden_dim=2048 in all four; head params exactly 18,188,632 (LLaVA) /
%    4,557,144 (CLIP); trained head is the cached HeadOnlyAgent, hidden width
%    held constant across backbones (NOT the 768-wide frozen_vision_baseline head).
%  * split_by_trajectory=true AND keep_best=true in all four. Cite THESE as the
%    reported 2x2; the legacy frame-level clip_combined_lang_stride4 must NOT be
%    cited (it would reintroduce the leakage Section~\ref{sec:split} avoids).
%  * Every optional knob off in all four: weighted_loss, learnable_bce_temp,
%    focal_gamma, past_action_slot_dropout, cam_weighted_loss, frame_history_k,
%    feature_norm/FiLM, image_dropout, patch_grid. (Head state-dict is exactly a
%    plain Linear->ReLU->Linear; CLIP feature_dim=1536=2x768 proves patch_grid=0.)
%  * Provenance (STEVE-1/VPT via MineDreamer play) confirmed.
%  * Hardware/stack: RTX 5090, torch 2.7 / cu128 (CUDA 12.8), bf16 + SDPA, no FA2.
%  * Loss (BCE + 0.5*(CE_x+CE_y)), Adam lr=1e-3 / batch 256 / 10 epochs / constant
%    LR / no wd-clip-warmup, fp16 cache -> fp32, model selection by val movement-F1,
%    decode (binary tau=1.0, camera tau=2.0 via eval_suite.sh), stride 4, four
%    combined-dataset caches: all confirmed.
%  * shah2021basalt correctly absent and uncited (data are VPT-agent rollouts).
%  * Python version corrected above to 3.12 (stack validated on 3.12; 3.10 is
%    unavailable on LMU CIP).
% ---------------------------------------------------------------------
% STILL OPEN:
%  1. Bibliography: dbstmpl.bib is still the stock LMU template. NONE of the 10
%     \cite keys here exist (Bain1995AFF, ross2011,
%     kumar2022finetuningdistortpretrainedfeatures,
%     li2023blip2bootstrappinglanguageimagepretraining,
%     liu2024improvedbaselinesvisualinstruction, radford2021clip, baker2022vpt,
%     tavakoli2018actionbranching, lifshitz2023steve1, kanervisto2020), and there
%     is no MineDreamer entry. The .bib must be populated before the chapter
%     compiles.
%  2. Random-policy baseline (see [VERIFY] in Evaluation Metrics): confirm a run
%     was logged, or keep the softened wording.
% ---------------------------------------------------------------------
% NOTE on the controllability protocol (not in this chapter): the registered envs
% are task-specific single biomes (MineRLChopATree640-v0 = forest,
% MineRLCollectDirt640-v0 = plains); no generic world with both trees and dirt
% exists in eval_envs.py and would need registering. A log breaks within the
% rollout budget under the *Fast* variants (BreakSpeedMultiplier=5.0, shared by
% both ...Fast-v0 and ...FastAim-v0); FastAim additionally sets start_pitch=20.0.
% =====================================================================